\documentclass[aps,prd,twocolumn,superscriptaddress,showpacs]{revtex4-2}
\usepackage{amsmath,amssymb,amsthm}
\usepackage{graphicx}
\usepackage{hyperref}
\hypersetup{colorlinks=true,linkcolor=blue,citecolor=blue,urlcolor=blue}

\newtheorem{theorem}{Theorem}
\newtheorem{lemma}{Lemma}

\begin{document}

\title{The Digital Standard Model:\\
Machine-Checked Structure of the Finite Spectral Triple}

\author{Jeffrey Michael Gurd}
\affiliation{Nexus Research (independent research)}
\email{jeffgurd888@gmail.com}

\begin{abstract}
We present machine-checked results on the finite spectral triple underlying
the Standard Model in noncommutative geometry. Using the Lean~4 proof
assistant, we classify the order-one commutant $W_{22}$ as a
22-dimensional real linear space---10 Standard Model directions plus 12
exotic survivor directions---with zero \texttt{sorry}s and no new axioms.
We further prove that the 12-dimensional exotic component decomposes
canonically as $8 \oplus 2 \oplus 2$: an 8-dimensional connected sector
and two 2-dimensional sectors, each generating a 3-dimensional compact
Lie algebra and commuting exactly with all directions outside itself.
We also machine-check two modular-dynamics results: the thermal modular
flux vanishes identically for the Gibbs-type Hamiltonian, and a
self-adjoint driver provably fails to commute with the off-diagonal
Dirac operator. All proofs are public and independently checkable. We
state precisely what is \emph{not} claimed: no physical particles, no
mass predictions, no continuum limit. This is the foundation of a
digital standard model: the Standard Model's finite geometry, made
machine-checkable, numerically explorable, and openly auditable.
\end{abstract}

\maketitle

\section{Introduction}

The Standard Model of particle physics can be formulated as a
noncommutative geometry: the finite spectral triple
$(\mathcal{A}_F, \mathcal{H}_F, D_F)$ with
$\mathcal{A}_F = \mathbb{C} \oplus \mathbb{H} \oplus M_3(\mathbb{C})$
and $\mathcal{H}_F = \mathbb{C}^{32}$, due to Chamseddine and
Connes~\cite{CC2007,CCM2007}. In this framework the gauge bosons,
the Higgs, and the fermion masses all emerge from the Dirac operator
and the spectral action. It is arguably the most economical known
mathematical encoding of the Standard Model's bosonic sector.

But economy is not the same as certainty. The finite triple is a
concrete object---explicit matrices, explicit commutators---and
concrete objects can be checked by machine. This paper reports what
happens when they are: a classification of the order-one deformations
of the finite Dirac operator, carried out in the Lean~4 proof
assistant~\cite{lean4} with every step verified and every gap closed.
We call the resulting program the \emph{digital standard model}:
not a new physical theory, but the existing finite geometry made
fully explicit, machine-checked, and numerically explorable.

Our contributions are:
\begin{enumerate}
\item A machine-checked classification (\texttt{cf\_kernel\_classification\_46\_22}):
the order-one commutant $W_{22}$ is a 22-dimensional real space,
spanned by 10 Standard Model directions and 12 exotic survivor
directions, each with explicit proved properties.
\item A machine-checked decomposition: the 12-dimensional exotic
component splits canonically as $8 \oplus 2 \oplus 2$.
\item Two machine-checked modular-dynamics theorems: thermal flux
vanishing and active-driver existence.
\item A public numerical laboratory corroborating and exploring
these structures.
\end{enumerate}
All Lean sources are public~\cite{thetlogos}; the classification
paper is archived~\cite{zenodo}.

A note on language that we enforce throughout: a 12-dimensional
\emph{algebraic} kernel is not 12 physical fields. The exotic
directions are explicit matrices with proved properties. They are
not claimed to be particles, forces, or dark matter. Interpretation
never alters an operator identity.

\section{The finite spectral triple}

We work with the established finite triple of
Chamseddine--Connes--Marcolli~\cite{CC2007,CCM2007}, recalled here
to fix notation. The algebra is
\begin{equation}
\mathcal{A}_F = \mathbb{C} \oplus \mathbb{H} \oplus M_3(\mathbb{C}),
\end{equation}
acting on $\mathcal{H}_F = \mathbb{C}^{32}$ via the representation
$\pi$, with opposite representation $\pi^\circ = J_F \pi(\cdot)^*
J_F^{-1}$. The finite Dirac operator $D_F$ is self-adjoint,
\[
D_F = \begin{pmatrix} 0 & M \\ M^* & 0 \end{pmatrix},
\qquad D_F^2 = \begin{pmatrix} MM^* & 0 \\ 0 & M^*M \end{pmatrix},
\]
with $M$ built from the Yukawa matrices
$Y_\nu, Y_e, Y_u, Y_d$ and the Majorana mass $Y_R$.
Charge conjugation $J_F$ satisfies $J_F^2 = I$, and the chiral
grading $\gamma_F$ satisfies $J_F \gamma_F = -\gamma_F J_F$.
The order-zero condition $[\pi(a), \pi^\circ(b)] = 0$ and the
order-one condition $[[D_F, \pi(a)], \pi^\circ(b)] = 0$ are the
axioms selecting admissible Dirac operators; both are
machine-checked in our formalization
(\texttt{order\_zero\_condition}, \texttt{OrderOneHolds\_smDirac}).

Nothing in this section is new; it is the common ground.

\section{The 22-dimensional classification}

\begin{theorem}[Machine-checked]
Let $W_{22}$ be the space of finite Dirac operators satisfying
the order-one condition with respect to the Standard Model
finite triple. Then $\dim_{\mathbb{R}} W_{22} = 22$, spanned by
the explicit list \texttt{dirs22}: 10 Standard Model directions
and 12 exotic survivor directions.
\end{theorem}

The proof (\texttt{ThetLogos.cf\_kernel\_classification\_46\_22})
proceeds by linear algebra over an explicit constraint system:
with \texttt{smGen : Fin 12}, the $12^2 = 144$ generator pairs
produce $144 \times 1024 = 147{,}456$ complex component equations,
and the kernel is computed exactly. The 10 SM directions are
5 complex moduli $\times$ 2 real
($y_\nu, y_e, y_u, y_d, y_R$). Each of the 12 exotic directions
carries four proved properties: self-adjointness,
$J$-compatibility, grading-oddness, and commutation with the
central element \texttt{cfMat} (48 lemmas in
\texttt{CFKernelRetarget.lean}). The formalization contains zero
\texttt{sorry}s; the axiom footprint is exactly
\texttt{propext}, \texttt{Classical.choice}, \texttt{Quot.sound},
as reported by \texttt{\#print axioms}.

The census runs in the 272-real-dimensional admissible Dirac
space. A separate 24-generator Python scan (576 pairs) corroborates
the classification numerically; it is supporting evidence only
and never a substitute for the machine proof.

\section{Exotic sector decomposition}

The 12 exotic directions are not an unstructured list. Define
the commutator graph on the 12 directions, with an edge when the
commutator is nonzero. Then:

\begin{theorem}[Machine-checked]
The 12-dimensional exotic component decomposes as
$8 \oplus 2 \oplus 2$:
\begin{enumerate}
\item an 8-dimensional sector (the NuLeR, ELNuR, NuREbarR,
EREbarR pairs) forming a connected commutator component;
\item a 2-dimensional sector (the ULDR pair) commuting exactly
with every direction outside itself;
\item a 2-dimensional sector (the DLUR pair) commuting exactly
with every direction outside itself.
\end{enumerate}
Each pair is internally non-abelian
($[\mathrm{Re}, \mathrm{Im}] \neq 0$, with the $(2,2)$ resp.\
$(3,3)$ entry equal to $-2i$); the ULDR pair spans with its
commutator a 3-dimensional space (proved
\texttt{LinearIndependent}); the DLUR analogue is the identical
proof shape.
\end{theorem}

The proof (\texttt{ThetLogos.ExoticSectors}, 78 theorems, zero
\texttt{sorry}s) works through all $2 \times 8 + 2 \times 8 +
2 \times 2 = 36$ cross-sector commutators, each shown to vanish
via disjoint-support lemmas, and exhibits the nonzero internal
commutators explicitly. The sector structure was first found by
numerical exploration of the explicit matrices and then
formalized; the machine proof is the result, the numerics were
the scout.

We emphasize: this is a decomposition of \emph{matrices}, not
of particles. The sector names reflect the fermion bilinears
the matrices couple, not physical identifications.

\section{Modular dynamics}

Let the modular flux be $\phi = i[K, D]$ for a modular
Hamiltonian $K$. Two results are machine-checked in
\texttt{ModularFlux.lean}:

\begin{theorem}[Machine-checked]
For the Gibbs-type modular Hamiltonian $K_\beta = \beta D^2$,
the modular flux vanishes identically:
$\phi = i[K_\beta, D] = 0$.
\end{theorem}

\begin{theorem}[Machine-checked]
For the off-diagonal Dirac operator
$D = \left(\begin{smallmatrix} 0 & M \\ M^* & 0 \end{smallmatrix}\right)$
with $M \neq 0$, there exists a self-adjoint operator $G$ with
$[G, D] \neq 0$ (witnessed explicitly by the block-diagonal
active driver).
\end{theorem}

The first promotes a thermal-stationarity observation to a
proved identity; the second closes an existence placeholder,
providing the non-commuting driver the dynamics requires.
Both are algebraic mechanisms. Neither proves anything about
physics.

We distinguish $K_\beta = \beta D^2$ (used above) from the
synthesized modular Hamiltonian $K_\rho = -\ln \rho$ of
\texttt{ModularTime.lean}; conflating them was an error in an
earlier draft, corrected here. The flux $\phi$ is always
lowercase: $\Phi$ is reserved as an ontological glyph and never
denotes an operator in this work.

A conditional spectral-gap lemma is also proved
(\texttt{spectralGap\_pos}): the gap $\Delta = \min$ positive
eigenvalue of a Hermitian $D$ is positive \emph{when $D$ is
invertible on its support}. The hypothesis is the content;
there is no unconditional mass-gap theorem, and the closed-form
gap is recorded as a blocked question (see Sec.~\ref{sec:not}).

\section{Numerical laboratory}

Alongside the proofs we maintain a deterministic numerical
companion suite (Tier-4 exploration, explicitly not proof):
a 24-generator census corroborating the classification, Connes
distances between states ($d_D = 0.4065\ \mathrm{GeV}^{-1}$ on
6 of 105 sampled pairs), thermal/KMS stationarity checks
(KMS error $6.77 \times 10^{-16}$ at $\beta = 1.0$), and an
interactive laboratory in which the 22 moduli can be varied
and the spectrum recomputed live. Numerical minima are never
presented as physical vacua; the laboratory is an instrument
for forming hypotheses, and every hypothesis it suggests must
earn its own proof.

\section{What is not claimed}
\label{sec:not}

To keep the record straight, we list what this work does
\emph{not} establish:
\begin{itemize}
\item No new particles, forces, or dark-matter candidates. The
exotic directions are matrices.
\item No mass predictions. Yukawa couplings are inputs; no
renormalization-group running is performed.
\item No mass gap. Only the conditional lemma above exists;
the closed-form spectral gap is a blocked question, obstructed
by the Galois-theoretic barrier on the degree-32 characteristic
polynomial and the non-analyticity of the $\min$ selection.
\item No MOND, no fifth forces, no cosmological predictions.
\item No continuum limit and no general-relativistic recovery.
\item No technology implementations.
\end{itemize}
Claims beyond this list that appeared in earlier drafts
(including a spacetime-dimension derivation from a withdrawn
kernel-dimension claim and a golden-ratio mass hierarchy) were
explicitly killed and are recorded as such in the public
proof-status ledger~\cite{thetlogos}. Killed claims return
only as new, weaker entries supported by fresh proof.

\section{Outlook}

Each rung of the program is conditional on the last:
decompose the 8-dimensional sector's internal structure;
generalize to $N$ generations; classify admissible deformations
$D_F \to D_F + \delta D$; connect the finite result to the
spectral action; and, finally, derive a quantitative, testable
prediction. The finite-to-continuum correspondence is a major
program, not an assumption.

\section{Conclusion}

The finite geometry underlying the Standard Model can be made
fully digital: explicit matrices, machine-checked theorems,
numerical laboratories, and an auditable record of what is
proved, what is explored, and what is dead. That is the
digital standard model---not a new theory of nature, but a new
standard of honesty about the mathematics we have. The proofs
are public; the reader is invited to check them.

\begin{thebibliography}{9}
\bibitem{CC2007} A.~Chamseddine and A.~Connes,
\textit{Why the Standard Model}, J.\ Geom.\ Phys.\ \textbf{58} (2008) 38.
\bibitem{CCM2007} A.~Chamseddine, A.~Connes, and M.~Marcolli,
\textit{Gravity and the Standard Model with neutrino mixing},
Adv.\ Theor.\ Math.\ Phys.\ \textbf{11} (2007) 991.
\bibitem{lean4} L.~de Moura \textit{et al.}, \textit{The Lean 4
Theorem Prover and Programming Language}, 2021.
\bibitem{thetlogos} J.~M.~Gurd, \textit{thet-logos},
\url{https://github.com/jeffgurd888/thet-logos}.
\bibitem{zenodo} J.~M.~Gurd, \textit{Thermal Time and Emergent
Space from Ontological Equality. Active Cartesian Modular
Dynamics (ACMD): the W22 Classification, Formalized in Lean},
Zenodo (2026), \doi{10.5281/zenodo.23140813}.
\end{thebibliography}

\end{document}
